\documentclass{situationist-preprint}

\title{On the Moral Weight of Farmed Animals}
\msid{ETH-10-b}
\msfamily{ETH-10}
\msdiscipline{Normative and Applied Ethics}
\mstype{Companion manuscript}
\mskeywords{moral weight; farmed animals; equal consideration; factory farming; capacity for welfare}

\begin{document}
\maketitle

\begin{abstract}
ETH-10-a argues that factory farming is impermissible and that humane farming is permissible but not morally neutral. A common reply is that farmed animals, having lower moral status than human beings, have interests of less weight, so that the pleasures of consumers may outweigh their suffering. We argue that the reply fails on the account of moral status in ETH-07. Under equal consideration, a given welfare change has the same weight whoever undergoes it; a lower capacity for welfare limits how much a being can lose, but the suffering of factory farming consumes a large part of what is at stake for the animals. The verdict on factory farming is therefore robust under every admissible assignment of status, while the reason against humane farming is real and modest.
\end{abstract}

\section{The discounting reply}

The argument of ETH-10-a does not weigh. Its principle of unnecessary suffering (Principle~P1 of ETH-10-a) forbids imposing severe suffering on beings with moral status for benefits that are not necessary, and it needs from ETH-07 only that the animals have moral status (Assumption~2.4 of ETH-10-a). A critic may nonetheless object that the principle is plausible only because we tacitly give the animals' suffering a weight comparable to human suffering. If chickens and pigs have much lower status than we do, the critic says, their suffering counts for much less, and the benefits to consumers, though not necessary, may outweigh it. This is the discounting reply.

\section{Weight without discount}

ETH-07-a distinguishes two questions: the degree of a being's status, and the weight of a given change in its welfare. By Theorem~3.4 of ETH-07-a, the order of beings by status is fixed by their capacity for welfare. By the principle of equal consideration (Definition~5.1 of ETH-07-a), a welfare change of a given magnitude has the same moral weight whichever being undergoes it, and the degree of status of a being is the total of what can be at stake for it. A lower status therefore means that less can be at stake for a being; it does not mean that what is at stake counts for less.

\begin{proposition}\label{p:weight}
Under equal consideration, the moral weight of the suffering imposed by a practice on an animal equals the magnitude of the welfare loss it imposes, and does not depend on the animal's status except through that magnitude.
\end{proposition}

\begin{proof}
Immediate from Definition~5.1 of ETH-07-a; compare Proposition~1.3 of ETH-07-c, which states the general form.
\end{proof}

The discounting reply therefore has to show that the welfare losses of factory farming are small. They are not. Broiler chickens are bred for growth so rapid that a large proportion develop leg disorders impairing their ability to walk, and they live with this condition for much of their short lives \citep{Knowles2008}. Breeding sows confined in stalls cannot turn around for most of a pregnancy. Whatever the capacity for welfare of a chicken or a pig, prolonged pain, frustration of the most basic behavioral needs and chronic distress lie at the bottom of the range open to it. A loss that takes a being to the bottom of its range is a large fraction of its capacity, however that capacity compares with ours. The animals concerned are mammals and birds, which MND-23 argues are conscious, so there is no doubt that the losses are experienced.

The number of animals multiplies the weight. Tens of billions of land animals are slaughtered each year, most of them chickens raised in industrial conditions. Against this the benefits that Assumption~2.5 of ETH-10-a identifies as unnecessary, namely pleasures of taste, convenience and lower prices that can be obtained or replaced by other means, are small per person and largely recoverable.

\begin{proposition}\label{p:robust}
The impermissibility of factory farming holds on every admissible assignment of moral status in the sense of Definition~2.7 of ETH-07-a, including hierarchical ones.
\end{proposition}

\begin{proof}
Every admissible assignment gives a sentient being positive status (Theorem~3.2 of ETH-07-a), and so satisfies Assumption~2.4 of ETH-10-a. Theorem~3.1 of ETH-10-a then follows from Principle~P1 and Assumption~2.5 of ETH-10-a, neither of which mentions the degree of status.
\end{proof}

Proposition~\ref{p:robust} shows that the verdict does not wait on the weighing. Proposition~\ref{p:weight} shows that, if one insists on weighing, the weighing does not favor the practice either: the hierarchical assignments that would discount animal suffering count capacity twice, as ETH-07-a argues in its Section~5, and even they leave severe suffering at a weight far above the unnecessary benefits.

\section{Humane farming}

The same account fixes the weight of the reason against humane farming. Its source is deprivation (Principle~P2 of ETH-10-a): killing an animal whose future would be good deprives it of that future. The weight of the reason is the magnitude of the loss, and Assumption~2.6 of ETH-10-a holds that the animal's time-relative interest in continuing to live is weak, in the sense of \citet{McMahan2002}: the animal is only weakly connected psychologically to its own future. The loss counted is therefore modest, and the weight of the reason is modest, though not zero. This is why humane farming is permissible and not morally neutral. It is a consequence of equal consideration, and not a discount on the animal's status: a human being with equally weak psychological connections to her future would suffer, in being killed, a loss of the same time-relative weight.

\section{Objection and reply}

\begin{objection}
Equal consideration ignores the difference that status makes. On hierarchical views, a given gain or loss of welfare matters more when it occurs in the life of a being with richer psychological capacities, such as a person, than in the life of an animal \citep{Kagan2019}.
\end{objection}

\begin{reply}
If richer capacities enlarge what can be at stake for a being, by giving it projects, relationships and a conception of its own future, equal consideration already counts them, since they enlarge the range of possible welfare changes. If it is supposed to enlarge the weight of a given pain beyond its magnitude, then two equal pains, one in a person and one in an animal, differ in weight although they do not differ as experiences. ETH-07-a rejects this as counting capacity twice. Nothing in the special standing of persons requires it.
\end{reply}

\section{Conclusion}

The moral weight of the suffering of farmed animals is the weight of the welfare it destroys, and the welfare it destroys is a large part of what those animals can have. The verdict of ETH-10-a on factory farming is robust under every admissible assignment of status, and the modest weight of the reason against humane farming reflects the modest loss that death imposes on beings weakly connected to their futures, not a lower weight for their interests.

\begin{thebibliography}{9}\small
\bibitem[Kagan(2019)]{Kagan2019} Kagan, S. (2019). \emph{How to Count Animals, More or Less}. Oxford University Press.
\bibitem[Knowles et~al.(2008)]{Knowles2008} Knowles, T.~G., Kestin, S.~C., Haslam, S.~M., Brown, S.~N., Green, L.~E., Butterworth, A., Pope, S.~J., Pfeiffer, D. and Nicol, C.~J. (2008). Leg disorders in broiler chickens: Prevalence, risk factors and prevention. \emph{PLoS ONE}, 3(2), e1545.
\bibitem[McMahan(2002)]{McMahan2002} McMahan, J. (2002). \emph{The Ethics of Killing: Problems at the Margins of Life}. Oxford University Press.
\end{thebibliography}

\end{document}
